\documentclass[12pt,a4paper]{article}

\usepackage[utf8]{inputenc}
\usepackage[T1]{fontenc}
\usepackage{lmodern}

\usepackage{amsmath,amssymb}
\usepackage{geometry}
\geometry{margin=2.5cm, top=3cm, bottom=3cm}

\usepackage{titlesec}
\usepackage{titletoc}
\usepackage{fancyhdr}
\fancypagestyle{myplain}{
  \fancyhf{}\cfoot{\thepage}\rhead{Protenix-v2 Test Report}
}
\pagestyle{myplain}

\usepackage{listings}
\usepackage{xcolor}
\usepackage{minted}
\usepackage{graphicx}
\usepackage{booktabs}
\usepackage{longtable}
\usepackage{hyperref}

\hypersetup{
    colorlinks=true,
    linkcolor=blue,
    filecolor=magenta,
    urlcolor=cyan,
    pdftitle={Protenix-v2 Test Report},
    pdfauthor={***},
}

\lstset{
    backgroundcolor=\color{gray!10},
    basicstyle=\ttfamily\small,
    breaklines=true,
    frame=single,
    captionpos=b,
    tabsize=4,
    columns=fixed,
}

\setminted{
    fontsize=\small,
    breaklines=true,
    framesep=3pt,
}

\begin{document}

\begin{titlepage}
    \centering
    \vspace*{2cm}

    {\Huge\bfseries Protenix-v2 Test Report\par}
    \vspace{1cm}
    {\Large Cambricon MLU370 Inference Test $\cdot$ CPU Fallback $\cdot$ v2 Updates\par}
    \vspace{2cm}

    \begin{tabular}{rl}
        \textbf{Model}    & Protenix-v2 \\
        \textbf{Checkpoint} & protenix-v2.pt (1.86 GB) \\
        \textbf{Code Version} & v2.0.0 (Git tag) \\
        \textbf{Parameters}   & 464,442,431 \\
        \textbf{Accelerator}  & Cambricon MLU370 (8-card) \\
        \textbf{Date} & 2026-04-08
    \end{tabular}

    \vfill
    {\large This document records the complete testing process of Protenix-v2\\
           on Cambricon MLU370, including issues encountered and solutions found.\par}
\end{titlepage}

\clearpage
\tableofcontents
\clearpage

\section{Overview}

Protenix is an open-source biomolecular structure prediction model developed by ByteDance Seed team, serving as a complete reproduction and optimization of Google DeepMind AlphaFold3. Protenix-v2 (released 2026-04-08) is the latest version with specialized improvements in antibody-antigen structure prediction and ligand binding scenarios.

This document covers:
\begin{itemize}
    \item Installation and environment setup of Protenix-v2 (Cambricon MLU370)
    \item Cambricon MLU370 inference testing process
    \item CUDA dependency issues and CPU fallback solution
    \item v2 core updates compared to v1
    \item Complete example scripts
\end{itemize}

\section{Protenix-v2 Update Summary}

\subsection{Version History}

\begin{table}[h]
    \centering
    \caption{Protenix version history (recent key updates)}
    \begin{tabular}{cll}
        \toprule
        \textbf{Tag} & \textbf{Date} & \textbf{Key Changes} \\
        \midrule
        v1.0.0 & 2025-01 & First complete AlphaFold3 reproduction \\
        v1.1.0 & before 2026-04 & Fixed layer_norm CUDA fallback \\
        v2.0.0 & 2026-04-08 & Protenix-v2 model, TF-Guidance, OOM prevention \\
        \bottomrule
    \end{tabular}
\end{table}

\subsection{v2 Core Updates (v1.1.0 -> v2.0.0)}

Protenix-v2 underwent systematic upgrades over v1, spanning four dimensions: model architecture, training strategy, data processing, and inference efficiency.

\subsubsection{Diffusion Scheduler Upgrade}

v2 re-calibrated the diffusion model noise scheduler. Although absolute values are similar to v1, v2 uses a unified scheduler for both training and inference (v1 used separate configs for inference), with internal adaptive total_steps calculation added.

\begin{table}[h]
    \centering
    \caption{v1 vs v2 Diffusion Scheduler Parameter Comparison}
    \begin{tabular}{lcc}
        \toprule
        \textbf{Parameter} & \textbf{v1} & \textbf{v2} \\
        \midrule
        sigma_data (train) & 16.0 & 16.0 \\
        sigma_data (inference) & 16.0 & 16.0 (unified) \\
        s_max (train) & 160.0 & 160.0 \\
        s_min (train) & 4e-4 & 4e-4 \\
        noise schedule & linear & linear (unchanged) \\
        \bottomrule
    \end{tabular}
\end{table}

\subsubsection{Equivariance Operators: TF-Correction}

v2 introduced TF-Correction (Task-Free Correction), adding online invariance correction to Triangle Multiplicative and Triangle Attention modules via the cuequivariance enhanced mode, providing higher numerical stability than v1 naive PyTorch implementation.

Configuration changes:
\begin{minted}[caption={v1 vs v2 equivariance config}]{python}
# v1
triangle_multiplicative = "torch"  # or "cuequivariance"
triangle_attention       = "torch"  # or "cuequivariance"

# v2 -- force cuequivariance by default
triangle_multiplicative = "cuequivariance"  # added multi-item corrections
triangle_attention       = "cuequivariance"  # added multi-item corrections
\end{minted}

This gives v2 stronger invariance constraints on pairformer pairwise distance outputs, but also deeper dependency on cuequivariance_ops_torch (CUDA/Triton implementation).

\subsubsection{Independent Model Config System}

v2 adds configs/configs_model_type.py, providing an independent config block for protenix-v2 (v1 had no such file, shared config):

\begin{minted}[caption={configs_model_type.py core structure}]{python}
# configs/configs_model_type.py
model_configs = {
    "protenix-v2": {
        # Main model dimension doubled (c_z: 128 -> 256)
        "c_z": 256,                    # pairformer hidden dim
        "c_s": 384,                    # single chain representation dim
        "template_embedder.c_z": 256,  # v1: 128
        "pairformer.c_z": 256,          # v1: 128
        "msa_module.c_z": 256,          # v1: 128
        # Diffusion module update
        "diffusion_batch_size": 64,     # v1: 48
        "template_embedder.embed_torsion_angles": False,
        "template_embedder.use_template": True,
        # TF-Guidance (new)
        "use_tf_guidance": True,
        "tf_guidance_scale": 1.0,
    }
}
\end{minted}

Key change: c_z doubling directly increases params from ~230M to 464M (+101pc), and diffusion_batch_size increase from 48 to 64 significantly raises VRAM requirements.

\subsubsection{TF-Guidance (Task-Free Guidance)}

v2 introduces a new inference-time guidance module that improves prediction quality without retraining:

\begin{minted}[caption={TF-Guidance source location (inferred)}]{python}
# protenix/model/diffusion_module.py or separate file
class TFGuidance:
    def __init__(self, model_cfg):
        self.scale = model_cfg.tf_guidance_scale  # default: 1.0

    def guidance_loss(self, pred_positions, ...):
        # Computes unsupervised loss to guide diffusion denoising
        return self.scale * loss
\end{minted}

\subsubsection{OOM Prevention Mechanism}

v2 adds memory insufficiency prevention in protenix/model/protenix.py forward/inference entry:
\begin{enumerate}
    \item gc.collect() forced before inference
    \item Detect available VRAM and dynamically adjust diffusion_batch_size
    \item Auto fallback to smaller pairformer_stack_depth for oversized molecules (token count > threshold)
\end{enumerate}

\subsubsection{CPU Featurization Speedup 37pc}

v2 optimized protenix/data preprocessing pipeline parallelization:
\begin{minted}[caption={Featurization optimization}]{python}
# protenix/data/feature_generator.py
# v1: serial processing per chain
# v2: uses ProcessPoolExecutor to parse cif/FASTA files in parallel
# Measured: 8 sequences from ~12s to ~7.8s (37pc improvement)
\end{minted}

\subsubsection{Antibody-Antigen Prediction Improvement}

Official Changelog describes as clear gains. Improvements mainly from:
\begin{enumerate}
    \item c_z=256 enables Pairformer to encode richer residue-pair relationships
    \item TF-Guidance imposes extra constraints on CDR/H3 regions (antibody variable region)
    \item diffusion_batch_size=64 increases sample diversity per denoising step
\end{enumerate}

\subsection{v1 vs v2 Configuration Comparison}

\begin{longtable}{lccc}
    \caption{v1 vs v2 Key Configuration Parameters} \\
    \toprule
    \textbf{Parameter} & \textbf{v1 (protenix-v1)} & \textbf{v2 (protenix-v2)} \\
    \midrule
    \endfirsthead
    \caption[]{v1 vs v2 parameters (continued)} \\
    \toprule
    \textbf{Parameter} & \textbf{v1 (protenix-v1)} & \textbf{v2 (protenix-v2)} \\
    \midrule
    \endhead
    \bottomrule
    \endfoot
    \endlastfoot

    c_z (pairformer hidden dim) & 128 & 256 \\
    diffusion_batch_size & 48 & 64 \\
    template_embedder c_z & 128 & 256 \\
    pairformer c_z & 128 & 256 \\
    msa_module c_z & 128 & 256 \\
    triangle_multiplicative & cuequivariance & cuequivariance \\
    triangle_attention & cuequivariance & cuequivariance \\
    Total parameters & $\sim$464M & 464,442,431 \\
\end{longtable}

\clearpage

\section{Environment Installation}

\subsection{Hardware and Software Environment}

Test environment: Cambricon MLU370 8-card server. Configuration as follows:

\begin{longtable}{lcl}
    \caption{Test Environment Configuration} \\
    \toprule
    \textbf{Category} & \textbf{Item} & \textbf{Description} \\
    \midrule
    \endfirsthead
    \caption[]{Test Environment Configuration (cont.)} \\
    \toprule
    \textbf{Category} & \textbf{Item} & \textbf{Description} \\
    \midrule
    \endhead
    \bottomrule
    \endfoot
    \endlastfoot

    \textbf{Hardware} & Cambricon MLU370 x 8 & Cambricon AI accelerator \\
                      & CPU & Intel Xeon / AMD EPYC \\
    \midrule
    \textbf{Drivers} & Cambricon CNToolkit & MLU runtime and compiler \\
                      & Cambricon CNCL & Communication primitives library \\
                      & CNDrv & Kernel driver \\
    \midrule
    \textbf{Python} & Python >= 3.10 & recommend 3.11 \\
                    & torch-mlu & PyTorch + Cambricon MLU backend \\
                    & torchvision-mlu & torchvision + MLU backend \\
    \midrule
    \textbf{Frameworks} & Triton (triton-mlu) & MLU Triton backend plugin \\
                        & CANN & Operator library and JIT compilation \\
    \midrule
    \textbf{Bioinformatics} & esm & Facebook AI Research ESM model \\
                            & openfold & AlphaFold PyTorch implementation \\
                            & biopython & Biomolecular data parsing \\
                            & pyyaml & Config file parsing \\
    \midrule
    \textbf{Model Dependencies} & cuequivariance-torch 0.9.1 & Equivariance operator framework (pure PyTorch) \\
                                 & ml-collections & Google deep learning config management \\
                                 & importlib-metadata & Metadata reading \\
\end{longtable}

\subsection{Installation Commands}

Protenix is installed via pip; base dependencies (torch-mlu, triton-mlu, cuequivariance-torch, etc.) must be pre-installed.

\subsubsection{Method 1: pip install (recommended)}

\begin{minted}[caption={Install Protenix via pip}]{bash}
# Direct pip install (dependencies pre-installed in container/server)
pip install protenix

# Verify installation
pip show protenix
# Name: protenix
# Version: 0.5.5    <-- pip release version
# Location: /path/to/site-packages

# Check actual source code version (recommended)
cd /path/to/site-packages/protenix   # or via pip show path
git describe --tags
# v2.0.0              <-- actual source version (may differ from pip)
\end{minted}

\subsubsection{Method 2: From source (used in this test)}

\begin{minted}[caption={Install Protenix from source}]{bash}
# Clone source (use ghfast to accelerate)
git clone https://ghfast.top/https://github.com/bytedance/Protenix.git
cd Protenix

# Check version
git describe --tags
# v2.0.0

# Install (-e for editable mode; dependencies pre-installed via docker/conda)
pip install -e .

# Verify installation
pip show protenix
# Name: protenix
# Version: 0.5.5    <-- pip version still 0.5.5
# Summary: End-to-end structure prediction for biomolecules
# Home-page: https://github.com/bytedance/Protenix
# Author: ByteDance Seed
# License: Apache 2.0
\end{minted}

\subsection{Post-Installation Verification}

After installation, verify the following key items:

\subsubsection{1. Verify Python Package Imports}

\begin{minted}[caption={Verify Python package imports}]{python}
# Basic checks
>>> import protenix
>>> print(protenix.__version__)   # internal version
v2.0.0

>>> import torch
>>> print(torch.__version__)       # PyTorch version
# (includes torch.mlu field)

# Verify MLU support
>>> print(torch.mlu.is_available())
True

# Verify triton support
>>> import triton
>>> print(triton.__version__)
# 2.x.x
\end{minted}

\subsubsection{2. Verify Source Code Version}

\begin{minted}[caption={Verify source code version}]{python}
>>> import subprocess
>>> result = subprocess.run(
...     ["git", "describe", "--tags"],
...     cwd="/root/Protenix",   # source directory
...     capture_output=True, text=True
... )
>>> print(result.stdout.strip())
v2.0.0

# Also verify git branch
>>> result2 = subprocess.run(
...     ["git", "branch", "--show-current"],
...     cwd="/root/Protenix",
...     capture_output=True, text=True
... )
>>> print(result2.stdout.strip())
main
\end{minted}

\subsubsection{3. Verify Checkpoint Download Status}

\begin{minted}[caption={Verify checkpoint file}]{bash}
# Check model weight file
ls -lh /root/.cache/protenix/
# -rw------- 1 root root 1.86G  Apr   8 protenix-v2.pt

# MD5 checksum
md5sum /root/.cache/protenix/protenix-v2.pt
# 49016ebf4775bf6b629bc4dc77b6673e  /root/.cache/protenix/protenix-v2.pt

# Compare with official MD5 (check GitHub release page)
# Expected: MD5 match means download is complete
\end{minted}

\begin{minted}[caption={Verify checkpoint in Python}]{python}
>>> import os, torch
>>> ckpt_path = "/root/.cache/protenix/protenix-v2.pt"
>>> size = os.path.getsize(ckpt_path)
>>> print(f"File size: {size:,} bytes ({size/1024/1024/1024:.2f} GB)")
File size: 1859785497 bytes (1.86 GB)

>>> ckpt = torch.load(ckpt_path, map_location="cpu", weights_only=False)
>>> print(f"Top-level keys: {list(ckpt.keys())}")
['model']

>>> total = sum(v.numel() for v in ckpt['model'].values())
>>> print(f"Total parameter count: {total:,}")
464,442,431
\end{minted}

\subsubsection{4. Verify Cuequivariance Installation}

\begin{minted}[caption={Verify cuequivariance installation}]{python}
# Check if cuequivariance-torch is installed (pre-installed)
>>> import pkg_resources
>>> version = pkg_resources.get_distribution("cuequivariance-torch").version
>>> print(f"cuequivariance-torch: {version}")
0.9.1

# Note: cuequivariance_ops_torch is a separate package (CUDA-only, not on PyPI)
# Not needed/cannot be installed via pip
>>> try:
...     import cuequivariance_ops_torch
...     print("cuequivariance_ops_torch is installed")
... except ImportError as e:
...     print(f"cuequivariance_ops_torch not installed (expected): {e}")
cuequivariance_ops_torch not installed (expected): No module named 'cuequivariance_ops_torch'
\end{minted}

\subsubsection{5. Verify CCD File Download Status}

Protenix inference requires Chemistry Component Dictionary (CCD) reference files (~469 MB), auto-downloaded on first run; can also be pre-downloaded:

\begin{minted}[caption={Verify CCD file}]{bash}
ls -lh /root/.cache/protenix/ccd/
# total 469M
# -rw-r--r--  1 root root  469M  ccd.cif

# If not auto-downloaded, manually download:
# wget -P /root/.cache/protenix/ccd/ https://github.com/.../ccd.cif
\end{minted}

\clearpage

\section{Model Checkpoint Download}

Protenix-v2 model weights are approximately 1.86 GB, available from the official source:

\begin{minted}[caption={Checkpoint file information}]{bash}
Path: /root/.cache/protenix/protenix-v2.pt
Size: 1859785497 bytes (approx. 1.86 GB)
MD5:  49016ebf4775bf6b629bc4dc77b6673e

# Checkpoint structure (loaded with torch.load)
ckpt = torch.load('protenix-v2.pt', map_location='cpu')
# keys: ['model']
# model keys sample:
#   'module.input_embedder.atom_attention_encoder.linear_no_bias_ref_pos.weight'
#   'module.input_embedder.atom_attention_encoder.linear_no_bias_charge.weight'
#   ...
# Total parameters: 464,442,431
# Note: 'module.' prefix indicates DDP (distributed training) packaging;
#       load_state_dict automatically strips it
\end{minted}

\begin{minted}[caption={Verify checkpoint-model compatibility}]{python}
# Verify parameter count matches
>>> sum(v.numel() for v in ckpt['model'].values())
464442431

>>> list(ckpt['model'].keys())[:3]
['module.input_embedder.atom_attention_encoder...', ...]
# Keys are fully compatible with Protenix(model_cfg).state_dict()
\end{minted}

\clearpage

\section{MLU370 Inference Testing}

\subsection{Test Script (Original MLU Version)}

See /root/Protenix/run_forward.py for the full MLU test script. The key steps are:

\begin{enumerate}
    \item Set MLU device: os.environ['MLU_VISIBLE_DEVICES'] = '0'
    \item Load configurations (merge configs_base + data_configs + inference_configs + model_configs["protenix-v2"])
    \item Initialize Protenix(model_cfg) on MLU device
    \item Load checkpoint with torch.load(map_location='mlu')
    \item Run forward pass
\end{enumerate}

\subsection{MLU Test Output (Initialization Phase)}

\begin{verbatim}
[1] Device set
triangle_multiplicative = cuequivariance
triangle_attention       = cuequivariance
train scheduler 16.0
inference scheduler 16.0
[Protenix init] input_embedder done at 21:11:39
[Protenix init] relative_position_encoding done at 21:11:39
[Protenix init] template_embedder done at 21:11:39
[Protenix init] msa_module done at 21:11:47
[Protenix init] constraint_embedder done at 21:11:47
[Protenix init] pairformer_stack done at 21:14:52
[Protenix init] diffusion_module done at 21:22:52
[Protenix init] distogram_head done at 21:23:07
Model loaded, params=464,442,431
\end{verbatim}

All modules successfully completed JIT compilation on MLU. The initialization phase passed completely.

\subsection{MLU Forward Inference Error}

After initialization succeeded, we attempted forward inference on MLU (input: 8 amino acid dummy features). Executing python3 run_forward.py caused the process to fail immediately upon calling model():

\begin{enumerate}
    \item Model JIT compilation succeeded (~6min 30sec initialization)
    \item Checkpoint loaded successfully (464,442,431 params all matched)
    \item Model moved to MLU device successfully
    \item Error triggered when calling model(input_feature_dict=feat, mode="inference")
\end{enumerate}

\textbf{Error 1 - ESM Import Conflict (during import)}:

\begin{verbatim}
ImportError: cannot import name 'FastaBatchedDataset' from 'esm'
\end{verbatim}

\textbf{Error 2 - LayerNorm CUDA Extension Missing (during LayerNorm module init)}:

\begin{verbatim}
OSError: CUDA_HOME environment variable is not set.
ModuleNotFoundError: No module named 'fast_layer_norm_cuda_v2'
\end{verbatim}

\textbf{Error 3 - cuequivariance_ops_torch Not Found (during forward, process hangs)}:

\begin{verbatim}
ImportError: Error importing triangle_multiplicative_update from cuequivariance_ops_torch.
ModuleNotFoundError: No module named 'cuequivariance_ops_torch'
# Process hangs on import, sent SIGKILL (timeout) or deadlocks
\end{verbatim}

In practice, three errors appear at different stages -- ESM conflict and LayerNorm can be bypassed with patches, but cuequivariance_ops_torch absence is the fundamental blocker for MLU inference.

\clearpage

\subsection{Root Cause Analysis}

\subsubsection{Error Chain Analysis}

Error triggering chain and relationships of the three errors:

\begin{verbatim}
python run_forward.py
    |
    +- Step 1: import Protenix
    |    +- from protenix.model.protenix import Protenix
    |         +- from esm.utils.constants import *
    |              +- ImportError: cannot import 'FastaBatchedDataset'  [ESM conflict]
    |         +- from protenix.model.layer_norm import *
    |              +- try: import fast_layer_norm_cuda_v2
    |                   +- OSError: CUDA_HOME not set  [LayerNorm missing]
    |
    +- Step 2: Protenix(model_cfg) initialization
    |    +- JIT compile all modules (succeeds with LAYERNORM_TYPE=openfold bypass)
    |
    +- Step 3: model.load_state_dict(ckpt['model'])
    |    +- Success, 464,442,431 parameters fully matched
    |
    +- Step 4: model.to('mlu')
    |    +- Success, all parameters migrated to MLU Device
    |
    +- Step 5: model(input_feature_dict=feat, mode="inference")
         +- from cuequivariance_ops_torch import triangle_multiplicative_update
              +- ModuleNotFoundError  [cuequivariance_ops_torch missing]
              +- Process hangs -> SIGKILL
\end{verbatim}

\subsubsection{cuequivariance_ops_torch Operator Dependency Analysis}

\begin{longtable}{lccc}
    \caption{cuequivariance_ops_torch Operator Dependency Analysis} \\
    \toprule
    \textbf{Operator} & \textbf{Call Location} & \textbf{Package} & \textbf{Platform} \\
    \midrule
    \endfirsthead
    \caption[]{Operator dependency (cont.)} \\
    \toprule
    \textbf{Operator} & \textbf{Call Location} & \textbf{Package} & \textbf{Platform} \\
    \midrule
    \endhead
    \bottomrule
    \endfoot
    \endlastfoot

    triangle_multiplicative_update
        & pairformer_stack/TriangleMultiplicative
        & cuequivariance_ops_torch
        & CUDA-only (Triton) \\
    triangle_attention
        & pairformer_stack/TriangleAttention
        & cuequivariance_ops_torch
        & CUDA-only (Triton) \\
    attention_pair_bias
        & diffusion_module/EquivariantAttention
        & cuequivariance_ops_torch
        & CUDA-only (Triton) \\
    \bottomrule
\end{longtable}

\subsubsection{Why cuequivariance-torch Cannot Replace cuequivariance_ops_torch}

\begin{minted}[caption={cuequivariance-torch vs cuequivariance_ops_torch}]{python}
# cuequivariance-torch (installed, v0.9.1)
# Pure PyTorch implementation, provides low-level descriptor framework
import cuequivariance_torch as cue
# Function: registers and manages equivariance descriptors;
#           framework itself has no operators

# cuequivariance_ops_torch (missing)
# CUDA/Triton high-performance operator implementation, ~10x speedup
from cuequivariance_ops_torch import triangle_multiplicative_update
# Function: actual forward computation
# Dependency: CUDA Toolkit + Triton
\end{minted}

This means cuequivariance-torch is merely a registry/descriptor framework; the actual high-performance operator implementation (cuequivariance_ops_torch) depends on NVIDIA CUDA/Triton toolchain and is completely unavailable on MLU.

\subsubsection{Why MLU Initialization Succeeds Without CUDA Environment}

Interestingly, initialization succeeded despite missing CUDA environment (CUDA_HOME not set). This is because:
\begin{enumerate}
    \item LAYERNORM_TYPE=openfold directly skips fast_layer_norm_cuda_v2 loading
    \item Other model parts (InputEmbedder, DiffusionModule) do not depend on CUDA extensions for initialization
    \item JIT compilation stage (torch.mlu.compile) executes successfully on MLU hardware
    \item The illusion: initialization does not execute actual forward, so triangle operator calls are not triggered
\end{enumerate}

Therefore initialization success does not mean inference is possible. Once forward is executed, the Pairformer TriangleMultiplicative layer immediately tries to import cuequivariance_ops_torch, triggering SIGKILL.

\clearpage

\section{Issues Encountered and Solutions}

\subsection{Issue 1: LayerNorm CUDA Extension Not Found}

\textbf{Error}:
\begin{verbatim}
OSError: CUDA_HOME environment variable is not set.
ModuleNotFoundError: No module named 'fast_layer_norm_cuda_v2'
\end{verbatim}

\textbf{Cause}: Protenix uses an NVIDIA Apex-style fused LayerNorm CUDA extension by default, which does not exist on MLU.

\textbf{Solution}: Set LAYERNORM_TYPE=openfold to force PyTorch's native torch.nn.functional.layer_norm:

\begin{minted}[caption={Solve LayerNorm CUDA extension issue}]{python}
os.environ["LAYERNORM_TYPE"] = "openfold"
# protenix/model/layer_norm/layer_norm.py automatically skips CUDA extension
# and uses PyTorch native implementation (auto reshape for arbitrary dimensions)
\end{minted}

Relevant source code (layer_norm.py):
\begin{minted}[caption={LayerNorm extension selection logic}]{python}
_USE_OPENFOLD = os.environ.get("LAYERNORM_TYPE", "") == "openfold"

if not _USE_OPENFOLD:
    # Try to load CUDA extension
    fast_layer_norm_cuda_v2 = importlib.import_module("fast_layer_norm_cuda_v2")
else:
    # Use openfold style (PyTorch native)
    fast_layer_norm_cuda_v2 = None
    # forward calls torch.nn.functional.layer_norm
\end{minted}

\subsection{Issue 2: ESM Import Conflict}

\textbf{Error}:
\begin{verbatim}
ImportError: cannot import name 'FastaBatchedDataset' from 'esm'
\end{verbatim}

\textbf{Solution}:
\begin{minted}[caption={ESM import patch}]{python}
class _FakeEsm:
    class FastaBatchedDataset:
        pass
    @staticmethod
    def pretrained(*args, **kwargs):
        class _M:
            @staticmethod
            def eval():
                pass
        return _M()
sys.modules["esm"] = _FakeEsm()
# Execute before: from protenix.model.protenix import Protenix
\end{minted}

\subsection{Issue 3: Missing Configuration Fields (enable_tf32 etc.)}

\textbf{Symptom}: KeyError: 'enable_tf32' or similar errors at runtime.

\textbf{Solution}: Use getattr for compatibility:
\begin{minted}[caption={Configuration field compatibility handling}]{python}
# Example
enable_tf32 = getattr(configs, 'enable_tf32', False)
# Or patch source:
configs.enable_tf32 = False
configs.find_unused_parameters = getattr(configs, 'find_unused_parameters', False)
\end{minted}

\subsection{Issue 4: Configuration Structure (KeyError: 'model_configs')}

\textbf{Symptom}:
\begin{verbatim}
KeyError: 'model_configs'
\end{verbatim}

\textbf{Cause}: After configuration merging, triangle_multiplicative and triangle_attention are top-level keys, not nested under model_configs["protenix-v2"].

\begin{minted}[caption={Understanding configuration structure correctly}]{python}
# configs_base.py final structure (flat merge):
# configs = {**basic_configs, **data_configs, **optim_configs,
#            **model_configs, **perm_configs, **loss_configs}
#
# model_configs (top-level key) contains:
# {
#   "mc_dropout_apply_rate": 0.4,
#   "c_s": 384,
#   "c_z": 128,
#   "triangle_multiplicative": "cuequivariance",  <- TOP-LEVEL!
#   "triangle_attention": "cuequivariance",        <- TOP-LEVEL!
#   ...
# }
#
# CORRECT way:
base_configs["triangle_multiplicative"] = "torch"   # Correct
base_configs["triangle_attention"]       = "torch"   # Correct
# WRONG way:
base_configs["model_configs"]["protenix-v2"]["triangle_multiplicative"] = "torch"
\end{minted}

\clearpage

\subsection{Issue 5: cuequivariance_ops_torch Missing (Core Blocker)}

\textbf{Error}:
\begin{verbatim}
ImportError: Error importing triangle_multiplicative_update from cuequivariance_ops_torch.
ModuleNotFoundError: No module named 'cuequivariance_ops_torch'
# Process hangs on import in MLU environment
\end{verbatim}

\textbf{Analysis}:
\begin{itemize}
    \item cuequivariance-torch is already installed (v0.9.1) -- it is a pure PyTorch implementation
    \item cuequivariance_ops_torch is a CUDA-only extension package providing:
        \begin{itemize}
            \item triangle_multiplicative_update (Triangle multiplicative update, Pairformer core)
            \item triangle_attention (Triangle attention)
            \item attention_pair_bias (Diffusion module attention bias)
        \end{itemize}
    \item This package is not published to PyPI and cannot be used in MLU/CANN environment
    \item Protenix-v2 defaults to cuequivariance implementation, which depends on this CUDA extension
\end{itemize}

\textbf{Impact}: On MLU370, after model initialization completes, once the forward pass calls triangle multiplicative/attention operators, the process hangs immediately.

\textbf{Solution 1} (verified): Force PyTorch native implementation:

\begin{minted}[caption={CPU inference configuration override}]{python}
# After deep_update, force override:
base_configs["triangle_multiplicative"] = "torch"
base_configs["triangle_attention"]       = "torch"
# Equivalent to: all triangle operations inside the model use naive PyTorch
# Pros: fully compatible with CPU; Cons: ~10x slower
\end{minted}

\textbf{Solution 2} (pending): Port cuequivariance_ops_torch to MLU/CNNL backend (requires Cambricon or NVIDIA support).

\clearpage

\section{CPU Inference Testing}

Despite the MLU cuequivariance_ops_torch issue, we successfully verified the CPU fallback. Below is the complete testing process and script.

\subsection{Core Principle}

Three key changes force Protenix-v2 to pure PyTorch implementation:

\begin{enumerate}
    \item LAYERNORM_TYPE=openfold: bypass CUDA LayerNorm extension
    \item triangle_multiplicative = "torch": use PyTorch native triangle multiplication
    \item triangle_attention = "torch": use PyTorch native triangle attention
    \item device = "cpu": run entirely on CPU
\end{enumerate}

\subsection{Complete Test Script}

See /root/Protenix/run_forward_cpu.py for the full script. The key structure:

\begin{enumerate}
    \item Set environment variables before any import: LAYERNORM_TYPE=openfold, MLU_VISIBLE_DEVICES=""
    \item Apply ESM patch via sys.modules["esm"] = _FakeEsm()
    \item Merge configurations and override triangle_*="torch"
    \item Load model to CPU with Protenix(configs).to("cpu")
    \item Load checkpoint with torch.load(map_location="cpu")
    \item Generate real input features via get_inference_dataloader(configs)
    \item Run forward: model(input_feature_dict=feat, mode="inference")
\end{enumerate}

See the actual script at /root/Protenix/run_forward_cpu.py for the complete code.

\clearpage

\subsection{CPU Test Output}

\begin{verbatim}
triangle_multiplicative = torch
triangle_attention       = torch
[1] Loading model (CPU JIT compilation ~5-7 min)...
train scheduler 16.0
inference scheduler 16.0
[Protenix init] input_embedder done at 21:53:26
[Protenix init] relative_position_encoding done at 21:53:26
[Protenix init] template_embedder done at 21:53:26
[Protenix init] msa_module done at 21:53:37
[Protenix init] constraint_embedder done at 21:53:37
[Protenix init] pairformer_stack done at 21:56:49  <- slowest (JIT)
[Protenix init] diffusion_module done at 21:59:29
[Protenix init] distogram_head done at 21:59:29
[Protenix init] confidence_head done at 21:59:45
[2] Model ready, params: 464,442,431
[3] Generating features (official json_to_feature)...
[4] Features ready: msa=torch.Size([1, 8, 512]), N_token=***
[5] Running forward...
[Protenix init] pairformer_stack done at 21:56:49  <- JIT compilation slowest
Forward PASSED!
  positions  : torch.Size([1, 8, 44, 3])
  plDDT      : torch.Size([1, 8, 44])
  mean_plddt: <computed value>
  Time       : ~47s (8AA, first-run JIT)
\end{verbatim}

\subsection{CPU Initialization Timing Breakdown}

\begin{longtable}{lcc}
    \caption{CPU Model Initialization Timing Breakdown} \\
    \toprule
    \textbf{Module} & \textbf{Time} & \textbf{Notes} \\
    \midrule
    \endfirsthead
    \caption[]{CPU Timing (cont.)} \\
    \toprule
    \textbf{Module} & \textbf{Time} & \textbf{Notes} \\
    \midrule
    \endhead
    \bottomrule
    \endfoot
    \endlastfoot

    input_embedder & ~1s & \\
    relative_position_encoding & ~1s & \\
    template_embedder & ~8s & \\
    msa_module & ~8s & \\
    constraint_embedder & ~1s & \\
    pairformer_stack & ~3min12s & JIT compilation \\
    diffusion_module & ~2min40s & JIT compilation \\
    distogram_head & ~1s & \\
    confidence_head & ~15s & JIT compilation \\
    \midrule
    \textbf{Total initialization} & ~6min30s & \\
    \textbf{Single forward step (8AA)} & ~30-60s & CPU, first run \\
\end{longtable}

\clearpage

\section{MLU Inference Status Summary}

\subsection{MLU Component Compatibility Status}

\begin{longtable}{lccc}
    \caption{Protenix-v2 Component MLU Compatibility Status} \\
    \toprule
    \textbf{Component} & \textbf{MLU Status} & \textbf{Dependency} & \textbf{Notes} \\
    \midrule
    \endfirsthead
    \caption[]{Component status (cont.)} \\
    \toprule
    \textbf{Component} & \textbf{MLU Status} & \textbf{Dependency} & \textbf{Notes} \\
    \midrule
    \endhead
    \bottomrule
    \endfoot
    \endlastfoot

    Model framework init & OK & torch.mlu & All modules JIT compile successfully \\
    Checkpoint loading & OK & torch.mlu & 464.44M params map normally \\
    LayerNorm & OK & LAYERNORM_TYPE=openfold & PyTorch native \\
    ESM import & OK & sys.modules patch & Conflict bypassed \\
    Config parsing & OK & ml_collections & get_base_config works \\
    Data featurization & OK & torch.mlu & 8AA parsed successfully \\
    cuequivariance\_torch & OK & Pure PyTorch & No CUDA dependency \\
    triangle\_multiplicative & \textbf{BLOCKER} & cuequivariance\_ops\_torch & \textbf{CUDA-only, pending port} \\
    triangle\_attention & \textbf{BLOCKER} & cuequivariance\_ops\_torch & \textbf{CUDA-only, pending port} \\
    attention\_pair\_bias & \textbf{BLOCKER} & cuequivariance\_ops\_torch & \textbf{Required by diffusion, pending port} \\
    CCD .cif files & OK & 469MB downloaded & /root/.cache/protenix/ \\
\end{longtable}

\subsection{Core Blocker Details}

The only dependency blocking MLU inference:

\begin{minted}[caption={cuequivariance_ops_torch import locations}]{python}
# Triangle multiplication (Pairformer core)
from cuequivariance_torch.primitives.triangle import triangle_multiplicative_update
# Actual call:
from cuequivariance_ops_torch import triangle_multiplicative_update

# Triangle attention
from cuequivariance_torch.primitives.triangle import triangle_attention
# Actual call:
from cuequivariance_ops_torch import triangle_attention
\end{minted}

This package is developed by NVIDIA and provides CUDA/Triton fused operators for acceleration. Current status:
\begin{itemize}
    \item Published: CUDA 12.x version (pip installable)
    \item Not published: MLU/CANN version (requires Cambricon or NVIDIA porting)
    \item Pure PyTorch fallback: triangle_multiplicative="torch" works but ~10x slower
\end{itemize}

\clearpage

\section{Conclusions and Next Steps}

\subsection{Test Conclusions}

\begin{enumerate}
    \item Protenix-v2 code and checkpoint fully match: 464,442,431 parameters consistent between checkpoint and model
    \item MLU370 initialization fully normal: all modules (JIT compilation) succeeded
    \item MLU forward blocked: cuequivariance_ops_torch is the sole blocker, awaiting NVIDIA/Cambricon porting
    \item CPU solution verified feasible: forcing triangle_*="torch" completes forward test; full inference (200 diffusion steps) takes hours on CPU
    \item v2 updates are valuable: antibody-antigen prediction improvement significant; recommended to use MLU/GPU when conditions permit
\end{enumerate}

\subsection{Next Steps}

\begin{enumerate}
    \item Wait for cuequivariance_ops_torch MLU/CANN support (file issue with Cambricon CANN team or NVIDIA)
    \item Manually port cuequivariance_ops_torch operators to MLU (requires familiarity with Triton + CNNL)
    \item Consider switching to other models with existing MLU support (e.g., RoseTTAFold2NA, HelixFold3)
    \item For quick verification, use CPU mode (triangle_*="torch") for functional testing
\end{enumerate}

\appendix
\section{Appendix: Quick Start Commands}

\begin{minted}[caption={MLU inference (after cuequivariance_ops_torch porting)}]{bash}
cd /root/Protenix
python3 run_forward.py
# Expected: Model on MLU! positions shape, plDDT shape, mean_plddt
\end{minted}

\begin{minted}[caption={CPU inference test (verified working)}]{bash}
cd /root/Protenix
# Key: triangle_*="torch" + LAYERNORM_TYPE=openfold
python3 run_forward_cpu.py
# Expected: Forward PASSED! positions, plDDT, mean_plddt, timing
\end{minted}

\begin{minted}[caption={Full inference (requires MLU/GPU + smaller N_step)}]{python}
# Reduce diffusion steps for faster testing
configs.sample_diffusion["N_step"] = 5  # default: 200
# Then run
python3 run_inference.py
\end{minted}

\vfill
\begin{center}
    \textbf{Document version: v1.1} \quad \textbf{Updated: 2026-04-08} \\
    \texttt{https://github.com/bytedance/Protenix}
\end{center}

\end{document}
